\subsection{Decision tree: how do I integrate this?}
\R{look at the integrand, take the first match}{
\begin{rcp}
\item \hd{In the table (\S9)?} \ra\ read it off. Inner function $ax+b$ \ra\ multiply by $\frac1a$. \warn nothing else may be "inner".
\item \hd{$\int f'(x)\,g(f(x))\dx$, i.e.\ the inner derivative is a factor?} \ra\ \textbf{substitution} $u=f(x)$.
Special case $\int\frac{f'}{f}\dx=\ln|f|+C$ and $\int f^n f'\dx=\frac{f^{n+1}}{n+1}$.
\item \hd{Product of two unrelated types} (poly $\times$ $e^x$/$\sin$/$\ln$/$\arctan$) \ra\ \textbf{by parts}.
\item \hd{Rational function $\frac{p}{q}$} \ra\ polynomial division if $\deg p\ge\deg q$, then \textbf{partial fractions}.
\item \hd{Rational in $\sin,\cos$} \ra\ \textbf{Weierstrass} $t=\tan\frac x2$; first try half-angle identities, often much shorter.
\item \hd{$\sqrt{a^2-x^2},\ \sqrt{a^2+x^2},\ \sqrt{x^2-a^2}$} \ra\ \textbf{trig substitution} $x=a\sin\theta$ / $a\tan\theta$ / $a\sec\theta$.
\item \hd{$\sqrt[n]{ax+b}$} \ra\ substitute $u=\sqrt[n]{ax+b}$ (kills the root completely).
\item \hd{$\sin^m x\cos^n x$} \ra\ odd power: peel one factor off, substitute the other. Both even: half-angle formulas.
\item \hd{Nothing works} \ra\ symmetry (\S6.9), or the answer contains a non-elementary integral like $\int\frac{\sin x}x\dx$ --- that is allowed in MC answers.
\end{rcp}
\hd{ALWAYS finish by differentiating your result.} That is the only real check under time pressure.
}

\subsection{Substitution \de{Substitution}}
\R{step by step (both directions)}{
\begin{rcp}
\item Choose $u=g(x)$ (usually the inner function of a composition).
\item Compute $du=g'(x)\dx$ \ra\ solve for $\dx=\frac{du}{g'(x)}$ and make \emph{every} $x$ disappear.
\item \hd{Definite integral:} convert the bounds $a\mapsto g(a)$, $b\mapsto g(b)$ --- then you must \emph{not} substitute back.
\item \hd{Indefinite integral:} integrate in $u$, then \textbf{substitute $u=g(x)$ back}.
\item \warn After back-substitution, differentiate: the chain rule factor $g'(x)$ must cancel. This is where most marks are lost.
\end{rcp}
\hd{Reverse direction} ($x=\varphi(t)$, needs $\varphi$ bijective/$C^1$): $\int f(x)\dx=\int f(\varphi(t))\varphi'(t)\dt$, back via $t=\varphi^{-1}(x)$.
}
\X{the chain-rule-back-substitution trap}{
$\int\frac{1}{2\cos^2\!\left(\frac x2\right)}\dx$: table gives $\int\frac{du}{\cos^2u}=\tan u$.
Substitute $u=\frac x2$, $\dx=2\,du$:
$\int\frac{2\,du}{2\cos^2u}=\tan u=\tan\!\frac x2+C$.
\textbf{NOT} $2\tan\frac x2$. The factor 2 from $\dx$ and the $\frac12$ in front cancel --- write the substitution out, never patch factors by feeling.
}

\subsection{Integration by parts \de{partielle Integration}}
\R{$\int u\,v'\dx=uv-\int u'v\dx$}{
\begin{rcp}
\item Choose $u$ = the part that gets \emph{simpler} when differentiated. Priority \hd{LIATE}: \textbf{L}og, \textbf{I}nverse trig, \textbf{A}lgebraic (poly), \textbf{T}rig, \textbf{E}xponential --- earlier in the list becomes $u$.
\item $v'$ = the rest; integrate it to get $v$ (no $+C$ needed here).
\item Apply the formula. Repeat if the new integral is still a product ($x^2e^x$: twice).
\item \hd{$\int\ln x\dx$, $\int\arctan x\dx$:} take $v'=1$. \ra\ $x\ln x-x$, $x\arctan x-\frac12\ln(1+x^2)$.
\item \hd{Cycle trick} ($\int e^x\sin x\dx$): after two rounds the original integral $I$ reappears \ra\ solve $I=\dots-I$ algebraically for $I$, divide by 2.
\end{rcp}
}
\E{multiple choice "$\int x\ln x\sin x\dx=\dots$"}{
Do \textbf{not} integrate --- \textbf{differentiate the four options} and compare with the integrand. 20 s instead of 5 min.\\
Check of $-x\ln x\cos x+\ln x\sin x+\sin x-\int\frac{\sin x}{x}\dx$:\\
$\frac{d}{dx}$: $-(\ln x+1)\cos x+x\ln x\sin x\ +\ \frac{\sin x}{x}+\ln x\cos x\ +\ \cos x\ -\ \frac{\sin x}{x}=x\ln x\sin x$ \checkmark
}

\subsection{Partial fractions \de{Partialbruchzerlegung}}
\R{$\int\frac{p(x)}{q(x)}\dx$}{
\begin{rcp}
\item $\deg p\ge\deg q$ \ra\ polynomial division first.
\item Factor $q$ completely over $\Rl$ (linear factors + irreducible quadratics).
\item Ansatz: $(x-a)$ \ra\ $\frac{A}{x-a}$;\ \ $(x-a)^k$ \ra\ $\frac{A_1}{x-a}+\dots+\frac{A_k}{(x-a)^k}$;\ \ irreducible $x^2+bx+c$ \ra\ $\frac{Ax+B}{x^2+bx+c}$ (and powers of it analogously).
\item Multiply out, compare coefficients \emph{or} plug in the roots (\hd{cover-up}: fastest for simple linear factors).
\item Integrate: $\int\frac{A}{x-a}=A\ln|x-a|$;\ \ $\int\frac{A}{(x-a)^k}=\frac{A}{(1-k)(x-a)^{k-1}}$;\ \
$\int\frac{Ax+B}{x^2+bx+c}$ \ra\ split into $\frac A2\cdot\frac{2x+b}{x^2+bx+c}$ (gives $\ln$) $+$ rest \ra\ complete the square \ra\ $\frac1a\arctan\frac{x-p}{a}$.
\end{rcp}
}

\E{four integrals, fully worked}{
\hd{(a) Substitution.} $\int\frac{x}{1+x^2}\dx$: $u=1+x^2$, $du=2x\dx$ \ra\ $\frac12\int\frac{du}{u}=\frac12\ln(1+x^2)+C$.\\
\hd{(b) By parts, twice.} $\int x^2e^x\dx$: $u=x^2,v'=e^x$ \ra\ $x^2e^x-2\int xe^x\dx$; again \ra\ $x^2e^x-2(xe^x-e^x)=e^x(x^2-2x+2)+C$.\\
\hd{(c) Partial fractions.} $\int\frac{1}{x^2-1}\dx=\int\left(\frac{1/2}{x-1}-\frac{1/2}{x+1}\right)\dx=\frac12\ln\left|\frac{x-1}{x+1}\right|+C$.\\
\hd{(d) Trig substitution.} $\int\sqrt{1-x^2}\dx$: $x=\sin\theta$, $\dx=\cos\theta\,d\theta$ \ra\ $\int\cos^2\theta\,d\theta=\frac\theta2+\frac{\sin2\theta}4$
\ra\ back: $\frac12\arcsin x+\frac12x\sqrt{1-x^2}+C$ (using $\sin2\theta=2x\sqrt{1-x^2}$). \warn back-substitution needs the identity, not just $\theta=\arcsin x$.
}

\subsection{Rational in $\sin/\cos$}
\R{half-angle first, Weierstrass only if needed}{
\hd{Try first:} $1+\cos x=2\cos^2\frac x2$,\quad $1-\cos x=2\sin^2\frac x2$,\quad $\sin x=2\sin\frac x2\cos\frac x2$.\\
\hd{Weierstrass} $t=\tan\frac x2$: $\ \sin x=\frac{2t}{1+t^2}$, $\cos x=\frac{1-t^2}{1+t^2}$, $\dx=\frac{2}{1+t^2}dt$ \ra\ rational in $t$ \ra\ \S6.4.
}
\E{$\int\frac{1}{1+\cos x}\dx$}{
$1+\cos x=2\cos^2\frac x2$ \ra\ $\int\frac{\dx}{2\cos^2\frac x2}
\overset{u=x/2}{=}\int\frac{2\,du}{2\cos^2u}=\tan u=\boxed{\tan\!\frac x2+C}$.\\
Check: $\frac{d}{dx}\tan\frac x2=\frac{1}{2\cos^2\frac x2}=\frac1{1+\cos x}$ \checkmark\\
\hd{Twin:} $\int\frac{\dx}{1-\cos x}=-\cot\frac x2+C$.
}

\subsection{Definite integrals, FTC, symmetry}
\F{Fundamental theorem (Hauptsatz)}{
$f$ continuous on $[a,b]$, $F(x):=\int_a^xf(t)\dt$ \ra\ $F$ differentiable and $F'=f$.\\
\hd{Newton--Leibniz:} $\int_a^bf=F(b)-F(a)$ for any antiderivative $F$.\\
\hd{Variable bounds (chain rule!):} $\frac{d}{dx}\int_{u(x)}^{v(x)}f(t)\dt=f(v(x))v'(x)-f(u(x))u'(x)$.\\
\hd{MVT for integrals:} $f$ cont.\ \ra\ $\exists\xi\in[a,b]:\int_a^bf=f(\xi)(b-a)$.
}
\F{free simplifications}{
\begin{bul}
\item $f$ \hd{odd}, symmetric interval \ra\ $\int_{-a}^{a}f=0$. $f$ \hd{even} \ra\ $=2\int_0^af$.
\item $\int_a^b=-\int_b^a$;\quad $\int_a^b=\int_a^c+\int_c^b$;\quad $\left|\int f\right|\le\int|f|$.
\item Monotonicity: $f\le g$ \ra\ $\int f\le\int g$. Standard estimate: $\left|\int_a^bf\right|\le(b-a)\max|f|$.
\item Periodic with period $T$: $\int_{a}^{a+T}f$ is independent of $a$.
\item $\int_0^{2\pi}\sin^2=\int_0^{2\pi}\cos^2=\pi$; $\int_0^\pi\sin x\dx=2$.
\end{bul}
}

\R{limit of a sum $=$ integral (Riemann sum trick)}{
If a limit looks like $\displaystyle\lim_{n\to\infty}\frac1n\sum_{k=1}^{n}f\!\left(\frac kn\right)$ then it equals $\displaystyle\int_0^1f(x)\dx$.
\begin{rcp}
\item Factor out $\frac1n$ (the width) --- if it is not there, the trick does not apply.
\item Write the rest as a function of $\frac kn$ only.
\item Read off $f$ and integrate over $[0,1]$. General interval: $\frac{b-a}{n}\sum f\!\left(a+k\frac{b-a}{n}\right)\to\int_a^bf$.
\end{rcp}
\hd{Ex.} $\lim\sum_{k=1}^n\frac{n}{n^2+k^2}=\lim\frac1n\sum\frac1{1+(k/n)^2}=\int_0^1\frac{\dx}{1+x^2}=\frac\pi4$.\\
\hd{Ex.} $\lim\frac1n\sum_{k=1}^n\ln\frac kn=\int_0^1\ln x\dx=-1$ (improper but convergent).
}
\R{integrals with $|\cdot|$, $\max$, or pieces}{
\begin{rcp}
\item Find the points where the sign / the active branch changes (zeros of the inside).
\item Split the integral there and drop the $|\cdot|$ with the correct sign on each piece.
\item Add up. \warn Never pull $|\cdot|$ out of the integral.
\end{rcp}
\hd{Ex.} $\int_0^2|x-1|\dx=\int_0^1(1-x)\dx+\int_1^2(x-1)\dx=\frac12+\frac12=1$.
}

\subsection{Riemann integrability \de{Riemann-integrierbar}}
\F{criteria (and the false friends)}{
\hd{Definition (Darboux):} $\sup_P L(f,P)=\inf_P U(f,P)$; equivalently $\forall\eps\exists P:U(f,P)-L(f,P)<\eps$.
\begin{bul}
\item \textbf{TRUE:} continuous on $[a,b]$ \ra\ integrable. Monotone on $[a,b]$ \ra\ integrable. Bounded with \emph{finitely many} (even countably many) discontinuities \ra\ integrable. Differentiable \ra\ continuous \ra\ integrable.
\item \textbf{TRUE (Lebesgue criterion):} bounded + discontinuities form a null set $\iff$ integrable.
\item \textbf{FALSE:} "must be continuous everywhere" --- jumps are fine.
\item \textbf{FALSE:} "integrable \ra\ differentiable" --- a step function is integrable, not even continuous.
\item \textbf{FALSE:} "finitely many function values \ra\ integrable" --- Dirichlet $1_\Q$ takes only the values $0,1$ and is \emph{not} integrable.
\item \warn integrability requires \textbf{bounded} $f$ (otherwise it is an improper integral).
\end{bul}
}

\subsection{Improper integrals \de{uneigentliche Integrale}}
\R{does $\int_a^\infty f$ or $\int_a^b f$ (pole) converge?}{
\begin{rcp}
\item \hd{Locate every problem point:} $\pm\infty$ and every point where $f$ is unbounded. \warn Split the integral so each piece has \emph{exactly one} problem; \emph{all} pieces must converge.
\item Write it as a limit: $\int_a^\infty f=\lim_{R\to\infty}\int_a^Rf$, $\int_a^bf=\lim_{\eps\to0^+}\int_a^{b-\eps}f$.
\item \hd{Benchmarks:} $\int_1^\infty\frac{\dx}{x^p}$ conv.\ $\iff p>1$;\quad $\int_0^1\frac{\dx}{x^p}$ conv.\ $\iff p<1$.
(Remember: at $\infty$ you need \emph{fast} decay, at $0$ you may only be \emph{mildly} singular.)
$\int_2^\infty\frac{\dx}{x\ln^px}$ conv.\ $\iff p>1$. $\int_0^\infty e^{-\lambda x}\dx=\frac1\lambda$ ($\lambda>0$).
\item \hd{Comparison:} $0\le f\le g$, $\int g<\infty$ \ra\ $\int f<\infty$; $f\ge g\ge0$, $\int g=\infty$ \ra\ $\int f=\infty$.
\item \hd{Limit comparison:} $\frac{f(x)}{g(x)}\to c\in(0,\infty)$ \ra\ both behave the same. Replace $f$ by its leading behaviour at the problem point.
\item \hd{Sign changes:} check absolute convergence first; if that fails, oscillating integrals like $\int_1^\infty\frac{\sin x}{x}$ still converge (Dirichlet/parts) --- \emph{conditionally}.
\end{rcp}
}
\E{which of these converge?}{
\begin{bul}
\item $\int_0^\infty\frac{\sin t}{t^3}\dt$: at $t\to0$, $\frac{\sin t}{t^3}\sim\frac1{t^2}$ \ra\ \textbf{diverges} (problem is at 0, not at $\infty$!).
\item $\int_0^\infty\frac{dt}{t^2+t^4}$: at $t\to0$, $\sim\frac1{t^2}$ \ra\ \textbf{diverges}.
\item $\int_1^\infty\frac{\cos^2t}{t}\dt$: $\cos^2t=\frac{1+\cos2t}{2}$ \ra\ contains $\int\frac{1}{2t}$ \ra\ \textbf{diverges}.
\item $\int_1^\infty\frac{\cos t}{t}\dt$: parts with $u=\frac1t$ \ra\ \textbf{converges} (conditionally).
\end{bul}
\hd{Moral:} always inspect \emph{both} ends before answering.
}
